\chapter{Introduction}

\begin{quote}
\begin{flushright}
{\sl Heere sweet Lord, at your service.}\\
Horatio, The Tragedie of Hamlet.
\end{flushright}
\end{quote}

This document covers the version \version{} of the library \name{},
and is itself on a beta development level.
Any suggestion should be sent to \mail.

The library design is also subject to modifications, so please refer
to this documentation if you face linkage problems after the upgrade
of the library.
You can also consult the file \verb|ChangeLog| at the root of the
package, as well as in each subdirectories to read more specific
changes.

\section{A brief overview}

\name{} is a collection of sublibraries, each of then aiming to cover
some specific needs.
These components are
\begin{description}
\item[C Tools]: collection of basic C utilities;
\item[Data strucutres]: basic data structures as lists and trees;
\item[Numerical utilities]: various basic numerical tools, as finite
  and central difference, special functions,\ldots;
\item[Statistical tools]: basic statistical analysis and distributions
  properties routines;
\item[Random numbers]: (pseudo-)random numbers generation;
\item[Quasi-Monte Carlo]: quasi-Monte Carlo sequences generation.
\end{description}
Each of these parts is described in more details in the next chapters.
A set of example files is also provided in the directory \texttt{examples}.

\section{Requirements and installation}
\label{sec:requirements}

\subsection{Environment variables}

In order to allow the various libraries to be found during the
compilation process and execution, please add
the following lines in the file $\verb|.bashrc|$, present at the
root of your home directory:
\begin{verbatim}
export LIBRARY_PATH=/usr/local/lib:$LIBRARY_PATH
export LD_LIBRARY_PATH=/usr/local/lib:$LD_LIBRARY_PATH
\end{verbatim}
This line can be placed at any position position in the
\verb|.bashrc| file, but be careful to use a raw text editor. From a
text console, your can reach your home directory by simply enter the
command
\begin{verbatim}
cd
\end{verbatim}

The place where the program and the libraries will be installed may vary
following the operating system.
If you benefit from administration rights, the installation paths will
be typically \verb|/usr/local/bin| for the programs and
\verb|/usr/local/lib| for the libraries.
If you have not these rights, the typical locations will be
\verb|$HOME/bin| and \verb|$HOME/lib|, where \verb|$HOME| is the user
repertory.
%For reference, we will design by \verb|INSTALL_PATH| the path where the programs are installed.

\subsection{OpenBLAS}
\label{sec:blas}

\oratio{} makes heavy use of \cblas{}, the C version of \blas{} (Basic Linear Algebra Subroutines) functions, that have to be installed first.
Any library implementing these routines should be compatible, but could possibly
require the modification of the file \verb|configure.in| to be
detected.
A recommended version is \textsf{OpenBLAS}, available at \url{http://www.openblas.net}.
It is based on \textsf{GotoBLAS2}, available at \url{http://www.tacc.utexas.edu/research-development/tacc-software/gotoblas2}, a \blas{} implementation proposed by the University of Texas at Austin.
Praised as one of the best implementations, it has been released under BSD license in 2010, but is no more maintained.

\textsf{OpenBLAS} installation requires to compile its source code. At the time of writing, the last version was 0.2.14. First, extract the code and compile it with the commands
\begin{verbatim}
tar zxvf OpenBLAS-0.2.14.tar.gz
cd OpenBLAS-0.2.14
make
\end{verbatim}
If you do not have the utility \textsf{wget} installed, you will also have to manually downoad LAPACK on Netlib.
The considered version of OpenBLAS requires LAPACK 3.5.0, available at the address \url{http://www.netlib.org/lapack/#_lapack_version_3_5_0}, but newer versions of OpenBLAS could rely on more recent LAPACK distributions.
The LAPACK archive has to be copied in the OpenBLAS main directory.

After the compilation, you can install the library with the command
\begin{verbatim}
sudo make install
\end{verbatim}
assuming that your account has enough administration rights.
By default, OpenBLAS is installed in the directory \verb|/opt/OpenBLAS|, but it is possible to specify the installation directory using the option \verb|PREFIX|, for instance
\begin{verbatim}
sudo make PREFIX=/usr/local install
\end{verbatim}
Since \verb|/usr/local| is the directory where \oratio{} is installed by defaults, it is a good idea to centralize the installation inside this directory.

\subsection{\oratio{} installation}

Once these dependencies are satisfied, the library can
be installed with the following instructions:
\begin{verbatim}
./configure
make
sudo make install
\end{verbatim}
Optimization compilation options can be specified as arguments of \verb|configure|, for instance:
\begin{verbatim}
./configure CFLAGS=''-O2'' FFLAGS="-O2"
\end{verbatim}

This documentation can also be generated using the command
\begin{verbatim}
make pdf
\end{verbatim}
which however requires a proper installation of \LaTeX.

\section{Conventions}

$x(i)$ designs the $i$-th component of the vector $x$.
When comparing two vectors $x$ and $y$ in $\rit^n$, an inequality $x \leq y$ is meant componentwise,
that is $x(i) \leq y(i),\ i = 1,\ldots,n$.

\section{Credits}

The development of the library has been initiated at FUNDP - the University of
Namur (Belgium), continued at the CERFACS (Toulouse, France), and is
currently pursued at the Universit\'e de Montr\'eal (QC, Canada).

Various people have contributed to this library.
Their names is written at the beginning of concerned files, and the
complete list is reproduced here, by alphabetical order: Adrien
Bonneu, Cinzia Cirillo, Pierre Derian, Jordan Goblet, Gabriel Schwanen.

\section{Miscellaneous}

\name{} stands for {\sl Operational Research, A Toolbox in Operation}.
Most of the routines were partially or totally written during
the development of the software {\sf AMLET} ({\sl Another Mixed Logit
Estimation Tool}).
The name is a reference to the Shapekespeare's masterpiece Hamlet.
Horatio was the best friend of Hamlet, as the library \name{} is the
primary companion of the software {\sf AMLET}.

The painting of the front cover is due to Eug\`ene Delacroix, in 1835,
under the name {\sl  Hamlet and Horatio in the Graveyard}.

